\chapter{Uvod}

\section{Informacioni sistemi}
\TODO{Hamurabijev zakonik, dodaj kao primere za državne i društvene procese}
Ideja radnog sistema postoji praktično onoliko dugo koliko postoji organizovana trgovina i proizvodnja. Na primer, još u Mesopotamiji (oko 3000. p.n.e.) postoji evidencija trgovine žitom, kreditnih ugovora, kamata i dugova na glinenim tablicama. Kasnije, u vreme Vavilona, Hamurabijev zakonik (oko 1750. p.n.e.) uspostavlja formalna pravila ne samo za poslovne odnose, već i za šire državne i društvene procese, uključujući regulisanje ugovora, odgovornosti, radnih obaveza i sankcija \cite{carmona2007accounting}.

\TODO{General System Theory/System Theory/Socio-technical Systems, svako definiše svoj ontološki minimum} 
\textbf{Radni sistem} \eng{Work system} se, prema Alteru \cite{alter2013work}, definiše kao organizovani sociotehnički sistem u kojem ljudi i/ili mašine obavljaju \textbf{procese} i \textbf{aktivnosti} koristeći informacije, tehnologiju i druge resurse kako bi proizveli određene proizvode ili usluge za jasno definisane interne ili eksterne korisnike. Osnovna ideja ovog pristupa jeste da radni sistem prikazuje \textbf{kako se posao u organizaciji zaista odvija u praksi}, nezavisno od stepena automatizacije ili nivoa tehnološke podrške.
 
Dalje, Alter u svom radu predstavlja osnovne elemente najopštijeg oblika radnog sistema. Istih devet elemenata primenjuju se na različite podtipove kao što su informacioni sistemi, projekti, lanci snabdevanja ili elektronska trgovina, dok udaljeniji tipovi sistema mogu, u zavisnosti od konteksta, dodavati ili izostavljati pojedine elemente:
\begin{enumerate}
\item \textbf{Procesi i aktivnosti} predstavljaju konkretan rad koji se obavlja unutar sistema rada. Oni mogu imati strukturu i biti strogo definisani nizovi koraka (procedura), ali mogu biti i nestrukturisani i zavisni od ljudskog zaključivanja. Ključno je da radni sistem ima makar jednu aktivnost, u suprotnom ne proizvodi nikakvu vrednost.
\item \textbf{Učesnici} su ljudi koji obavljaju rad u sistemu rada, bez obzira na to da li koriste informacione tehnologije ili ne. U mnogim slučajevima i sami kupci mogu biti učesnici sistema rada.
\item \textbf{Informacije} obuhvataju sve oblike podataka koje radni sistem kreira, koristi, skladišti, prenosi, ažurira ili briše tokom aktivnosti. Ovo ne obuhvata samo digitalne podatke u informacionim tehnologijama već i verbalne dogovore, dokumenta, implicitno znanje, beleške i sl.
\item \textbf{Tehnologije} obuhvataju fizičke i/ili digitalne alate koji pomažu učesnicima da obave posao ili da ga u potpunosti automatizuju. Tehnologije mogu biti od običnog čekića, preko računara i softvera, sve do potpuno automatizovanih digitalnih agenata.
\item \textbf{Proizvod} u sistemu rada predstavlja rezultat njegovog postojanja, jer radni sistem kao krajnji izlaz proizvodi određenu informaciju, fizičku stvar i/ili uslugu koja donosi vrednost korisniku. Procena uspešnosti sistema rada može se meriti na osnovu proizvoda.
\item \textbf{Kupci} su oni koji dobijaju i koriste proizvode ili usluge sistema rada. Oni mogu biti spoljašnji (kupci organizacije) ili unutrašnji (npr. drugi odsek u kompaniji).
\item \textbf{Okruženje} predstavlja spoljne faktore koji utiču na rad sistema: zakone i regulativu, tržišnu konkurenciju, kulturu, standarde, organizacione politike, istoriju firme i sl. Ignorisanje okruženja može dovesti do pogrešne analize i neuspeha pri dizajnu sistema.
\item \textbf{Infrastruktura} obuhvata resurse na koje se radni sistem oslanja i koristi, ali koji nisu direktno pod njegovom kontrolom. To mogu biti drugi sistemi, ljudska infrastruktura, zajedničke baze podataka, cloud platforme i sl.
\item \textbf{Strategije} određuju pravac i nameru. One postoje na nivou preduzeća, odeljenja i konkretnog sistema rada. Strategije moraju biti usklađene i koherentne, jer neusaglašene strategije dovode do neuspešne implementacije i pogrešnih prioriteta.
\end{enumerate}

Prethodna definicija radnog sistema može se posmatrati kao analitička, jer pruža opis njegove strukture i međusobnih odnosa komponenti. Da bi se obezbedilo dublje teorijsko razumevanje pojma, uvodi se \textit{ontološka definicija} koja razmatra minimalni skup elemenata neophodnih da bi radni sistem postojao kao realan entitet. Iako Alter eksplicitno ne koristi termin ,,ontološki minimum radnog sistema’’, iz njegove definicije i logike pristupa proizlazi da su za postojanje svakog radnog sistema neophodni sledeći elementi:
\begin{enumerate}
\item \textbf{procesi i aktivnosti} - koji određuju \textit{šta se radi}
\item \textbf{učesnici} - koji određuju \textit{ko radi}
\item \textbf{resursi} (informacije i tehnologije) - koji određuju \textit{čime se radi}, i
\item \textbf{proizvod/ishod} - koji određuje \textit{šta se postiže}.
\end{enumerate}
Ovi elementi čine jezgro ontološkog okvira svakog radnog sistema i predstavljaju njegovu minimalnu strukturu nezavisno od konteksta ili stepena složenosti.

Radni sistem omogućava sveobuhvatno \textbf{analiziranje, razumevanje, evaluaciju i unapređenje} poslovnih aktivnosti, jer posmatra međusobno usklađen odnos svojih komponenti. Prema Alteru, radni sistemi postoje da bi stvarali vrednost za svoje korisnike, pa se njihova uspešnost ocenjuje ne samo prema efikasnosti unutrašnjih procesa, već i prema percepciji i zadovoljstvu korisnika proizvedenim proizvodima i uslugama.

U zavisnosti od svrhe i konteksta, pojam \textit{vrednosti} može se interpretirati na različite načine, što otvara prostor za teorijsko proširivanje Alterove definicije radnog sistema. Polazeći od različitih konteksta u kojima radni sistemi funkcionišu, u daljem radu oni se posmatraju kroz četiri osnovne kategorije: poslovne, javne, neprofitne i lične radne sisteme.

Poslovni sistemi teže stvaranju ekonomske vrednosti za klijente i organizaciju, dok su javni sistemi usmereni na stvaranje društvene vrednosti kroz pružanje usluga građanima i unapređenje javnog interesa \cite{moore1997creating}. Neprofitni sistemi fokusirani su na stvaranje zajedničkog dobra bez cilja ostvarivanja profita \cite{murray1994evaluating}, dok lični radni sistemi opisuju načine na koje pojedinci koriste tehnologiju i informacije za postizanje ličnih ciljeva i produktivnosti.

Bez obzira na kategoriju, svi radni sistemi dele zajednički ontološki minimum, procese i aktivnosti koje obavljaju učesnici koristeći informacije i tehnologije radi stvaranja određene vrednosti za jasno definisane korisnike. Kada se radni sistemi integrišu sa informacionim tehnologijama \cite{turban2000information}, koje obuhvataju hardver, softver, baze podataka i mreže kao primarna sredstva za procesiranje, čuvanje, razmenu i prikaz informacija neophodnih za realizaciju poslovnih procesa, tada govorimo o \textbf{Informacionim sistemima} kao posebnoj kategoriji radnih sistema. U daljem radu fokus će biti upravo na tim sistemima.